\section{The image of the leading projection}
\label{sec:image}

Throughout this section the coefficient field is $\F$.
Recall that $\overline L_n$ denotes the reduction of the saturated
$\Z_{(2)}$-lattice of rational solutions of weight $n$, and that
$F^r$ consists of polynomials containing at least $r$ occurrences of $B$
in the free associative algebra on
\[
 A=x^2,\qquad C=[x,y],\qquad B=y.
\]
We use the ordinary free Lie algebra on $A,C,B$ inside this associative
algebra, as justified in Lemma~\ref{lem:abc-free}.
For $\psi\in F^r\overline L_n$, let $\chi(A,C,B)$ be its component
with exactly $r$ occurrences of $B$, and put
\[
 f(C,B)=\chi(0,C,B).
\]
If this component is nonzero, then $n=2m+r$, where $m$ is the total
number of $A$'s and $C$'s in $\chi$.
Proposition~\ref{prop:projection-injective} says that this projection
is injective on the corresponding associated graded space.
We shall identify a free Lie algebra containing its image.

The argument begins with a restriction on products of iterated adjoints,
obtained from symmetry in $x,y$ in Lemma~\ref{lem:even-indices}.
This restriction gives invariance under the polynomial shifts
$B\mapsto B+s(C)$, and a Vandermonde argument then places $f$ in the
ordinary free Lie algebra generated by $\ad_C^kB$, $k\geq1$.
This proves the all-weight dimension bound.  We then establish a
separate compatibility: projected leading Ihara brackets become
ordinary brackets.  It will provide the independent Hall words in
Section~\ref{sec:induction}.

\subsection{Symmetry and polynomial shifts}

\begin{lemma}[The even-index bound]
\label{lem:even-indices}
Suppose that $n>2$, $\psi\in F^r\overline L_n$, and $n=2m+r$.
The component of $\psi$ of original $(x,y)$-bidegree $(m,m+r)$
is an associative linear combination of products
\begin{equation}
 \label{eq:swapped-products}
 (\ad_y^{j_1}x)\cdots(\ad_y^{j_m}x)
\end{equation}
in which at least $r$ of the indices $j_1,\ldots,j_m$ are even,
with zero counted as even.
Its image under $A=0$ is exactly $f(C,B)$.
In particular, if $r>m$, then $f=0$.
\end{lemma}

\begin{proof}
By the elimination argument in Lemma~\ref{lem:abc-free}, every
homogeneous Lie polynomial in $x,y$ of weight greater than one is a
Lie polynomial in $U_j=\ad_x^j y$, $j\geq0$.
In characteristic two these elimination letters are
\[
 U_{2i}=E_i=\ad_A^i B,
 \qquad
 U_{2i+1}=O_i=\ad_A^i C.
\]
Associative words in the $E_i,O_i$ are linearly independent.
Indeed, order the letters with $A$ largest and use degree followed
by lexicographic order.  The leading words of these letters are
$A^iB,A^iC$, respectively, with coefficient one.
They form a prefix code: one reads a string of $A$'s up to its
next $B$ or $C$, thereby determining each successive codeword.
Distinct products consequently have distinct leading words.

Each $E_i$ is homogeneous of $B$-count one, whereas each $O_i$
has $B$-count zero.  Thus the expansion of $\psi$ in independent
associative words in this alphabet is graded by the number of
$E$-factors.  The hypothesis $\psi\in F^r$ forces the coefficients
of all words with fewer than $r$ such factors to vanish separately.
This establishes a termwise restriction before any substitution
or quotient is taken.

Select the component of original bidegree $(m+r,m)$.
Every $U_j$ has exactly one $y$, so every product in this component
has exactly $m$ factors.  The reduced antisymmetry equation gives
$\psi(x,y)=\psi(y,x)$.  Swapping $x$ and $y$ therefore gives
the desired expression~\eqref{eq:swapped-products} for bidegree
$(m,m+r)$, with the same lower bound on the number of even indices.
If $r>m$, no such product exists.

For completeness, an $A,C,B$-word with respective counts $(a,c,b)$
has original bidegree $(2a+c,c+b)$.  In bidegree $(m,m+r)$ this means
\begin{equation}
 \label{eq:bidegree-counts}
 c=m-2a,\qquad b=r+2a.
\end{equation}
The terms with $a>0$ have higher $B$-count and vanish under $A=0$;
the terms with $a=0$ have $c=m,b=r$ and are precisely $f$.
Thus the even-index bound applies to the whole bidegree component,
including every term subsequently removed by the projection.
\end{proof}

\begin{lemma}[Polynomial shift invariance]
\label{lem:shift-invariance}
For every $f$ obtained as above and every polynomial $s(C)$,
\begin{equation}
 \label{eq:polynomial-shift}
 f(C,B+s(C))=f(C,B).
\end{equation}
The identity holds after adjoining any finite collection of
commuting scalar indeterminates to the coefficients; in particular,
it is a polynomial identity in the coefficients of $s$.
\end{lemma}

\begin{proof}
Let $R$ be any commutative $\F$-algebra, and work in
$S=R\langle C,B\rangle$.  Define an $R$-linear derivation $D$ by
\[
 DB=C,\qquad DC=0.
\]
In characteristic two the square of a derivation is again a
derivation: the two middle terms in the twice-applied product rule
cancel.  Since $D^2$ vanishes on $B,C$, it vanishes on all of $S$.
Consequently the Ore extension, followed by the indicated quotient,
\[
 \mathcal E=S[X;D]/(X^2),\qquad Xq=qX+Dq\quad(q\in S),
\]
has underlying left $S$-module $S\oplus SX$.
Indeed $X^2q=qX^2+D^2q=qX^2$, so $X^2$ is central in the Ore
extension and its quotient introduces no relation in $S$.
In particular, $S$ embeds in $\mathcal E$.
The substitutions $x\mapsto X$, $y\mapsto B$ give
\[
 x^2\mapsto0,\qquad [x,y]\mapsto C.
\]
Equivalently, $\mathcal E$ acts on $S$ with $X$ acting by $D$
and elements of $S$ acting by left multiplication.
For an element written $q+q'X$, its value on $1$ is its constant
term $q$.

Fix $s\in R[C]$, and put $F=B+s(C)$.
The substitution $B\mapsto F$, $C\mapsto C$ commutes with $D$,
since $D(s)=0$.  It therefore extends to an automorphism of
$\mathcal E$ fixing $X$.
Let $h(x,y)$ be the whole bidegree-$(m,m+r)$ component in
Lemma~\ref{lem:even-indices}.  Evaluation at $x=X,y=B$ kills the
terms containing $A$ in its $A,C,B$ expression, so that
$h(X,B)=f(C,B)$ in the embedded subalgebra $S$ of $\mathcal E$.
Applying the shift automorphism just constructed gives
$h(X,F)=f(C,F)$.  We may therefore evaluate the even-index product
expansion of $h$ term by term in $\mathcal E$.  Its sum belongs to
$S$, so its value on $1$ is exactly $f(C,F)$.
The factors of positive index become elements of $S$:
\[
 \ad_F^jX=P_j:=\ad_F^{j-1}C\quad(j\geq1),
\]
whereas the index-zero factor is $X$.
We shall bound the $B$-degree of the constant term of each product
separately.

Write
\[
 L=\ad_B,\qquad Q=\ad_{s(C)},\qquad K=\ad_C,
 \qquad P_j=(L+Q)^{j-1}C.
\]
Here $L$ raises $B$-degree by one and $Q$ preserves it.
The relations needed below are
\[
 Q(C)=0,\qquad QK=KQ,\qquad DL=LD+K,\qquad DQ=QD.
\]
For the two smallest positive indices one has
\begin{equation}
 \label{eq:small-indices}
 P_1=C,\quad DP_1=0,\qquad
 P_2=[B,C],\quad DP_2=0.
\end{equation}
For $j\geq2$, every nonzero operator word in $(L+Q)^{j-1}C$
has rightmost operator $L$, because $Q(C)=0$.
Thus $P_j$ has no $B$-degree-zero component, and its degree-one
component is
\[
 (P_j)_{[1]}=Q^{j-2}LC.
\]
This is killed by $D$, since $DLC=KC=[C,C]=0$.
As $D$ lowers $B$-degree by one, $DP_j$ has no degree-zero
component for any $j\geq1$.

For $j\geq3$, the degree-two component is explicitly
\[
 (P_j)_{[2]}=
 \sum_{a+b=j-3}Q^aLQ^bLC.
\]
Upon applying $D$, differentiating the rightmost $L$ gives zero,
and differentiating the other $L$ gives $Q^aKQ^bLC$.
The commutation of $Q$ and $K$ therefore gives
\begin{equation}
 \label{eq:differentiated-linear-part}
 (DP_j)_{[1]}=(j-2)Q^{j-3}KLC.
\end{equation}
For even $j$ this is zero in $R$.
Together with~\eqref{eq:small-indices}, these calculations prove
the following lower bounds; the zero polynomial satisfies every
listed bound:
\[
\begin{array}{c|cc}
 & P_j & DP_j\\ \hline
 j\geq1\text{ odd} & 0 & 1\\
 j\geq2\text{ even} & 1 & 2
\end{array}
\qquad\text{(minimum $B$-degree).}
\]
The bound for odd $P_j$ is deliberately weak, so that it includes
$P_1=C$.
All these calculations allow constant terms in $s$ and arbitrary
scalar coefficients in $R$.

Now consider one product with $q$ index-zero factors and $e$ positive
even-index factors.  Move each $X$ to the right using
$XP=PX+DP$, and apply the resulting expression to $1$.
Every surviving term must use each of the original $q$ copies of
$X$ as a differentiation, since a terminal $X$ kills $1$.
Repeated use of the product rule assigns these differentiations
to positive-index factors to their right.  An allocation assigning
two differentiations to the same factor vanishes, because $D^2=0$.
Hence every surviving term is a product in which exactly $q$
distinct positive-index factors have been differentiated once.
An undifferentiated positive even-index factor contributes a lower
bound of one, and an odd-index factor contributes a lower bound
of zero.  Differentiating either type increases the respective
lower bound by one.  The resulting $B$-degree is therefore at least
\[
 e+q,
\]
which is precisely the number of all even indices, including zero.
If no such allocation is possible the constant term is zero, which
also satisfies the bound.  Cancellations among allocations cannot
create lower-degree terms.

Lemma~\ref{lem:even-indices} now shows that $f(C,B+s(C))$ has no
$B$-degree below $r$.  On the other hand, $f$ is homogeneous of
$B$-degree $r$, and replacing $B$ by $B+s(C)$ cannot increase that
degree.  Its degree-$r$ component is exactly $f(C,B)$, obtained by
choosing $B$ at every occurrence.  This proves
\eqref{eq:polynomial-shift} over $R$, and in particular over any
polynomial ring in scalar indeterminates.
\end{proof}

\subsection{The ordinary Lie image and the dimension bound}
\begin{proposition}[The image bound]
\label{prop:image-bound}
Set $g_k=\ad_C^kB$ for $k\geq1$.
These elements freely generate an ordinary free Lie algebra inside
$\Lie_{\F}\langle C,B\rangle$.
For $n>2$, the image of the leading projection of
$\operatorname{gr}^r_F\overline L_n$ belongs to its weight-$n$,
length-$r$ component, where $\wt(g_k)=2k+1$ and each $g_k$ has
length one.
The associated graded piece with $r=0$ is zero.
\end{proposition}

\begin{proof}
First suppose $r\geq1$.
Encode the associative words of $B$-degree $r$ by the vector-space
isomorphism
\[
 \mathcal P_r:
 C^{i_0}BC^{i_1}\cdots BC^{i_r}
 \longmapsto s_0^{i_0}s_1^{i_1}\cdots s_r^{i_r}
 \quad\text{in }\F[s_0,\ldots,s_r].
\]
Put $P=\mathcal P_r(f)$.
For $1\leq i\leq r$, let $R_iP$ be its restriction obtained by
identifying $s_{i-1}$ and $s_i$, written in the same polynomial
ring $\F[t_0,\ldots,t_{r-1}]$ for every $i$:
\[
 R_iP=P(t_0,\ldots,t_{i-2},t_{i-1},t_{i-1},
               t_i,\ldots,t_{r-1}).
\]
Empty initial or final strings in this formula are omitted.
In the coefficient of $\lambda$ in $f(C,B+\lambda C^j)$, replacing
the $i$th $B$ merges its two neighboring powers of $C$ and adds
$j$ to their exponent.  Its encoding with $r-1$ remaining $B$'s
is consequently $t_{i-1}^{j}R_iP$.
Lemma~\ref{lem:shift-invariance}, applied with the independent
scalar $\lambda$, gives
\begin{equation}
 \label{eq:vandermonde-system}
 \sum_{i=1}^r t_{i-1}^{j}R_iP=0\qquad(j\geq0).
\end{equation}
Taking $j=0,\ldots,r-1$ produces a Vandermonde matrix whose
determinant is
\[
 \prod_{0\leq a<b\leq r-1}(t_b-t_a).
\]
It is a nonzero polynomial over $\F$, because the $t_i$ are
independent indeterminates.  Thus the matrix is invertible over
$\F(t_0,\ldots,t_{r-1})$, and all $R_iP$ vanish.
This argument also applies to even $r$; it never divides by $r$
or specializes the $t_i$ to elements of the two-element field.
For $r=1$ the matrix is simply $(1)$.

The kernel of the $i$th identification is the principal ideal
generated by $d_i=s_{i-1}-s_i$.
These $d_i$ are distinct nonassociate prime linear polynomials.
Their individual divisibility therefore implies
\begin{equation}
 \label{eq:consecutive-divisibility}
 d_1\cdots d_r\mid P
 \quad\text{in }\F[s_0,\ldots,s_r].
\end{equation}
It remains to translate this associative statement into an
ordinary Lie statement.

Temporarily include $g_0=B$.
The ordinary Lie ideal generated by $B$ in
$\Lie_{\F}\langle C,B\rangle$ is the subalgebra generated by
$g_k=\ad_C^kB$, $k\geq0$.
Indeed, this subalgebra contains $B$ and is stable under $\ad_C$
by the derivation rule and $[C,g_k]=g_{k+1}$; the reverse inclusion
is immediate.  Since $f$ has positive $B$-degree, it belongs to
this ordinary Lie subalgebra.

Expanding the iterated adjoints gives the useful product formula
\begin{equation}
 \label{eq:gap-monomials}
 \mathcal P_r(g_{k_1}\cdots g_{k_r})
 =\prod_{i=1}^r(s_{i-1}-s_i)^{k_i}
 =d_1^{k_1}\cdots d_r^{k_r}.
\end{equation}
The variables $d_1,\ldots,d_r,s_r$ are an invertible linear
change of coordinates from $s_0,\ldots,s_r$.
In particular, the monomials in the right side of
\eqref{eq:gap-monomials} are linearly independent.
For different product lengths the $B$-degrees differ, so all
associative words in the $g_k$, $k\geq0$, are independent.
Their ordinary Lie algebra is therefore the ordinary free Lie
algebra on these letters, by its embedding in the free associative
algebra.

Write the associative expansion of $f$ uniquely in these independent
words.  Its encoding lies in $\F[d_1,\ldots,d_r]$.
By~\eqref{eq:consecutive-divisibility}, every monomial appearing
with nonzero coefficient has positive exponent of every $d_i$.
Equivalently, the expansion of $f$ contains no word with any factor
$g_0$.  The ordinary free Lie algebra on $g_0,g_1,\ldots$ has a
retraction
\[
 \rho(g_0)=0,\qquad \rho(g_k)=g_k\quad(k\geq1).
\]
Its extension to the free associative algebra fixes the expansion
of $f$.  Hence $\rho(f)=f$, proving that $f$ belongs to the
ordinary free Lie algebra on $g_k$, $k\geq1$.
This retraction preserves ordinary Lie membership throughout;
no identification with all primitive elements in characteristic
two is being made.

Each $g_k$ contains one $B$ and $k$ copies of $C$.
The original weight is therefore $2k+1$, and the $B$-count of a
Lie word is its length in these generators.  This proves the
asserted weight and length statement.

If $r=0$, then $f$ is an ordinary Lie polynomial in $C$ alone.
The ordinary Lie algebra on one generator is its one-dimensional
linear span, of original weight two; thus $f=0$ for $n>2$.
Proposition~\ref{prop:projection-injective} then gives
$\operatorname{gr}^0_F\overline L_n=0$.
The possible boundary $m=0$ also contributes nothing: its leading
component would be an ordinary Lie polynomial in $B$ alone,
whose only possible positive weight is one.
\end{proof}

Define the ordinary free Lie algebra
\[
 \mathfrak f_{\mathrm{odd}}
   =\Lie_{\F}\langle z_k:k\geq1\rangle,
 \qquad \wt(z_k)=2k+1,
\]
and give each $z_k$ length one.
Write $\mathfrak f_{\mathrm{odd},n,r}$ for its component of
weight $n$ and length $r$, and set
\begin{equation}
 \label{eq:free-dimensions}
 d_{n,r}=\dim_{\F}\mathfrak f_{\mathrm{odd},n,r},
 \qquad d_n=\sum_{r\geq0}d_{n,r}.
\end{equation}
These are also the corresponding dimensions over $\Q$, by
the ordinary free Lie lattice of Lemma~\ref{lem:lie-lattice}.
The sum is finite: a length-$r$ word has weight at least $3r$,
and only finitely many generators can occur in any fixed weight.

\begin{corollary}[The all-weight upper bound]
\label{cor:dimension-bound}
For every $n$ and $r\geq0$ there is an injection
\[
 \operatorname{gr}^r_F\overline L_n
 \lhook\joinrel\longrightarrow
 \mathfrak f_{\mathrm{odd},n,r},
\]
obtained by the leading projection and the identification
$g_k\leftrightarrow z_k$.  In particular,
\begin{equation}
 \label{eq:dimension-upper-bound}
 \dim_{\Q}W_n
 =\dim_{\F}\overline L_n
 \leq\sum_{r\geq0}d_{n,r}=d_n.
\end{equation}
\end{corollary}

\begin{proof}
For $n>2$, combine Proposition~\ref{prop:projection-injective}
with Proposition~\ref{prop:image-bound}.
If $n-r$ is odd or negative the source piece is zero by the
weights of $A,C,B$, and the target piece is likewise zero since
the weight of a length-$r$ word in odd-weight letters has parity
$r$.  The cases $r=0$ and $m=0$ were disposed of in
Proposition~\ref{prop:image-bound}.
Weights one and two have zero solution space by
Lemma~\ref{lem:low-weights}, and have no generators or Lie words
in $\mathfrak f_{\mathrm{odd}}$.
Finally, $F^{n+1}\overline L_n=0$, so this is a finite filtration
in weight $n$ and the dimensions of its associated graded pieces
sum to $\dim_{\F}\overline L_n$.
The equality of this dimension with $\dim_{\Q}W_n$ is
Proposition~\ref{prop:saturated-reduction}.
\end{proof}

\subsection{Leading Ihara brackets}
The upper bound concerns the full equation space.  To attain it in
Section~\ref{sec:induction}, we will construct elements whose projected
leading terms are the $g_k$ and compare their iterated Ihara brackets
with ordinary Lie words.  The following compatibility is computed in
the ambient Lie algebra and does not require closure of $W$.

\begin{proposition}[The leading Ihara bracket]
\label{prop:leading-ihara}
Let $\psi\in F^r\overline L_n$ and
$\phi\in F^s\overline L_{n'}$, where $n,n'>2$ and $r,s\geq1$.
Write $f,g$ for their leading projections at counts $r,s$.
Their Ihara bracket, computed in the ambient Lie algebra in
$x,y$, belongs to $F^{r+s}$ in the $A,C,B$ presentation, and
the projection of its component of count $r+s$ is
\begin{equation}
 \label{eq:leading-ihara}
 [f,g].
\end{equation}
This assertion does not require prior closure of the full solution
space under the Ihara bracket.
\end{proposition}

\begin{proof}
We first prove the filtration assertion before applying $A=0$.
The derivation $\delta=\ad_x$ acts on the free $A,C,B$ algebra by
\[
 \delta A=0,\qquad \delta B=C,\qquad \delta C=[A,B].
\]
Write $\delta=\delta_-+\delta_+$, where
\[
 \delta_-B=C,\quad \delta_-A=\delta_-C=0,
 \qquad
 \delta_+C=[A,B],\quad \delta_+A=\delta_+B=0.
\]
These two derivations respectively lower and raise the $B$-count
by one.  In particular, $\delta(F^r)\subset F^{r-1}$.
Directly from the definition of the Ihara derivation,
\begin{equation}
 \label{eq:ihara-abc-action}
 D_\psi A=0,\qquad
 D_\psi B=[B,\psi],\qquad
 D_\psi C=[C,\psi]+[B,\delta\psi].
\end{equation}
The last formula is the derivation identity applied to $[x,B]$.
Thus $D_\psi B\in F^{r+1}$ and $D_\psi C\in F^r$:
the outside $B$ in $[B,\delta\psi]$ restores the one count that
$\delta_-$ can remove.  Replacing a $B$ in a word of count $s$
adds at least $r$ to its count, and replacing a $C$ also adds at
least $r$; the derivative of an $A$ vanishes.
It follows that
\[
 D_\psi(F^s)\subset F^{r+s}.
\]
Interchanging $\psi,\phi$ and including their ordinary bracket
proves the asserted containment of the Ihara bracket.

Let $\chi,\eta$ be the components of $\psi,\phi$ of counts
$r,s$.  In~\eqref{eq:ihara-abc-action}, only $\chi$ can contribute
to the part of the derivation that raises count by exactly $r$.
Its action is
\[
 A\longmapsto0,\qquad
 B\longmapsto[B,\chi],\qquad
 C\longmapsto[C,\chi]+[B,\delta_-\chi].
\]
Indeed, a higher-count component of $\psi$ still has higher count
after applying $[B,\delta_-(-)]$, while
$[B,\delta_+\chi]$ raises the count by two more.
This derivation preserves the ideal generated by $A$, so it
descends under $A=0$.
The projection commutes with $\delta_-$, whose induced action
on $\F\langle C,B\rangle$ is the derivation $D$ from
Lemma~\ref{lem:shift-invariance}.
The induced leading action is consequently
\[
 B\longmapsto[B,f],\qquad
 C\longmapsto[C,f]+[B,Df].
\]
By Proposition~\ref{prop:image-bound}, $f$ is an ordinary Lie
polynomial in $g_k$, $k\geq1$; each of these is killed by $D$,
because
\[
 Dg_k=\ad_C^k(C)=0\quad(k\geq1).
\]
Thus the extra term $[B,Df]$ vanishes, and the induced derivation
on $\F\langle C,B\rangle$ is $q\mapsto[q,f]$.
One can also read $Df=0$ directly as the coefficient of $\lambda$
in~\eqref{eq:polynomial-shift} with $s(C)=\lambda C$.

The two Ihara derivations therefore project to $[g,f]$ and
$[f,g]$, respectively, and the ordinary bracket projects to
$[f,g]$.  Their specified combination is
\[
 [g,f]-[f,g]+[f,g]=[g,f]=[f,g]
\]
over $\F$.
All higher components of $\psi$ or $\phi$ have already been
excluded by the filtration calculation, so this proves
\eqref{eq:leading-ihara}.
\end{proof}
